\section{Introducción}
En una sesión práctica de física, el tiempo disponible debe permitir formular preguntas, ejecutar análisis e interpretar resultados. Sin embargo, la diversidad de sistemas operativos, bibliotecas y permisos puede convertir el inicio de la clase en una serie de instalaciones individuales. El problema no se limita a comodidad técnica: si cada estudiante parte de un entorno distinto, resulta más difícil distinguir un error conceptual de uno causado por la configuración.

Las buenas prácticas de computación científica relacionan organización de archivos, control de versiones y automatización con la posibilidad de comprender y repetir el trabajo \cite{wilson}. Incorporarlas a la formación permite que el estudiante aprenda tanto el contenido de un análisis como las condiciones que lo hacen verificable. No se trata de ocultar la infraestructura indefinidamente, sino de decidir qué complejidad abordar en cada momento de una secuencia docente.

En América Latina, las propuestas de computación transversal han planteado compartir protocolos, servicios y capacidades entre proyectos e instituciones \cite{transversal,whitepaper}. El término transversal se utiliza aquí para describir prácticas reutilizables más allá de un experimento: organizar código y datos, definir dependencias, ejecutar tareas, registrar resultados y comunicar cambios. La formación en física ofrece un contexto concreto para aprenderlas, mientras su utilidad puede extenderse a otras áreas científicas.

Este trabajo desarrolla el resumen aceptado como contribución 667 de ICHEP 2026 \cite{abstract}. Describe una arquitectura docente, sitúa su aplicación en SILAFAE XV y CSC América Latina 2026, y propone cómo evaluar sus resultados sin confundir rapidez de preparación con aprendizaje. Añade un demostrador editorial: los artículos de este proyecto se compilan mediante un flujo que conserva las fuentes y actualiza los PDF en el mismo repositorio. El alcance es un informe técnico y una propuesta metodológica; no se presenta un ensayo educativo controlado.
